\section{Discussion}
\label{sec:discussion}

The standard texts on infrared interpretation describe the region below about 1400~\wn{} as complex and rarely diagnostic on its own, and the carbonyl and X-H regions as the place where group frequencies are reliable~\cite{colthup1990introduction,bellamy1975infrared}. That description is correct, and this study gives it numbers: above 1500~\wn{} the median rule multiplies the odds of its group by three, and the best ones by ten or forty; below 1350~\wn{} the median rule multiplies them by one. What the description does not give, because the counting had not been done at this scale, is how uniform the failure is (26 of 38 fingerprint rules below 1.2 in Chemotion), that the exceptions are exactly the rules that ask for more than a position, and that a rule which is useless for confirming a group can still be good for excluding it. The absence of the 690~\wn{} ring bend, of the secondary amide N-H stretch or of the terminal alkyne C-H stretch each lower the odds of their group by a factor of six to seventeen, a direction the tables do not list. The expert systems that once encoded these tables did attach a confidence to each rule~\cite{woodruff1980pairs,andreev1996expirs}, but it was the author's estimate; the values reported here are measured, and they can be re-measured on any other collection.

The tables could not have said it, and the reason is structural. A correlation table is built from spectra of compounds that have the group: it records where the band is and roughly how it looks. It has no denominator. Sensitivity can be estimated from such a collection, specificity cannot, and it is specificity that the fingerprint lacks. A second reason is that the rules were written for a reader looking at a whole spectrum, who weighs each band against its neighbours, notices patterns and discards what does not fit; nobody applies the C-O rule for secondary alcohols in isolation. We tested each rule on its own, which is what a student does with a table on the first day and what any program does that encodes the table literally. The likelihood ratios reported here are therefore a fair measure of the rules as written, a lower bound for what an experienced reader extracts from the same bands, and a warning to anyone who would automate the table one row at a time.

Table~\ref{tab:failures} groups the rules that carry no evidence in either collection by what fills their window. Eleven are single-bond C-O and C-N stretches between 1000 and 1350~\wn{}; their false positives are dominated by aryl esters, aryl ethers and aromatic ketones, whose C-O and aryl-C(=O) stretches sit in the same windows with odds ratios of 10 to 44. Six are bending and wagging modes between 600 and 860~\wn{}, filled by the C-H wags of monosubstituted and para rings. Two are fingerprint rules whose ``strong'' qualifier removes true positives faster than false ones. Beyond these, three failure modes recur in rules that do carry evidence: intensity qualifiers on bands that are not in fact strong (Section~\ref{sec:intensity}), windows narrower than the spread of the band across the group (the nitrile windows catch half of the NIST nitriles; the terminal alkyne C$\equiv$C window misses the conjugated alkynes of Chemotion), and absence conditions that a second ring in the molecule violates.

\begin{table}[t]
\centering\small
\caption{Failure modes of the frozen rules. The first three rows cover the 19 group rules with \LRp{} below 1.2 in both collections; the last three recur in rules that carry evidence. Odds ratios are those of the groups most over-represented among the false positives (NIST, primary threshold).}
\label{tab:failures}
\begin{tabular}{@{}p{3.6cm}p{4.6cm}p{6.6cm}@{}}
\toprule
Failure mode & Rules affected & What the data show \\
\midrule
Single-bond C-O and C-N stretches, 1000 to 1350~\wn{} & 11 rules: C-O of primary, secondary and tertiary alcohols, phenols, saturated, aryl alkyl and diaryl ethers, acids; C-N of primary and secondary amines and secondary amides & Sensitivity 0.9 to 1.0, specificity 0.01 to 0.16. Almost every spectrum has a band in the window; false positives are dominated by aryl esters (OR 11 to 44), aryl ethers (8 to 15) and aromatic ketones (8 to 44), whose own C-O and aryl-C(=O) stretches fall there. \\
Bending and wagging modes, 600 to 860~\wn{} & 6 rules: alcohol O-H wag, primary and secondary amine N-H wags, secondary amide N-H wag, para-disubstituted benzene wag, trisubstituted alkene wag & Filled by the aromatic C-H wags: monosubstituted rings (OR 11 to 43) and para rings (OR 5 to 64) among the false positives. In Chemotion, with 93\% aromatic compounds, specificity 0.02 to 0.16. \\
``Strong'' qualifier that reverses the sign & 2 rules: saturated ketone C-C-C (1100 to 1230), saturated ester O-C-C (1030 to 1100) & The band is present in 85 to 99\% of the group by position but reaches the strong class in 18 to 42\%; the qualifier removes true positives faster than false ones and \LRp{} falls to 0.6 to 1.2. \\
\midrule
Intensity qualifier on a band that is not strong & Nitrile C$\equiv$N (both classes), acid O-H stretch, secondary amide N-H bend, ester and ketone skeletal stretches & Specificity rises to 0.85 to 1.00 but sensitivity falls to 0.03 to 0.44; \LRp{} looks excellent while the rule misses most of the group and \LRn{} approaches 1 (Sections~\ref{sec:intensity} and \ref{sec:nitrile}). \\
Window narrower than the spread of the band & Nitrile windows (20~\wn{}); terminal alkyne C$\equiv$C 2100 to 2140 & By position alone the nitrile windows catch 45 to 47\% of NIST nitriles; the Chemotion alkynes, mostly conjugated, sit above the C$\equiv$C window (\LRp{} 2.7 against 12 in NIST). \\
Absence condition violated by a second ring & Ortho and para benzene rules (no band at 680 to 700~\wn{}) & 74\% of Chemotion ortho compounds show the forbidden band, against 56\% of the others, because the molecule carries a monosubstituted phenyl elsewhere; \LRp{} 0.82. \\
\bottomrule
\end{tabular}
\end{table}

On the intensity words we make one concrete suggestion. ``Strong'', ``medium'' and ``weak'' in the tables describe how a band looks to the eye in a transmission spectrum of a small molecule, and as such they are useful. They are not thresholds, and the nitrile shows what happens when they are read as thresholds. A rule that is meant to be applied literally, by a student or a program, should state what it means: either a relative height (the ester carbonyl, when it falls in its window, is above 0.6 of the strongest band in 80 to 90\% of esters; the nitrile stretch is above 0.2 in 43 to 65\% of nitriles and above 0.6 in one in ten) or a word that describes the right property, and for the nitrile that word is ``sharp'' or ``isolated'', not ``intense''. Table~\ref{tab:rules} and the released result tables give the relative heights that would let the qualifiers be rewritten this way for every rule in the set.

The table is meant to be used the way likelihood ratios are used in diagnosis. A reader who suspects an ester from the reaction and sees a strong band at 1740~\wn{} multiplies the odds by 43 (12 in a collection rich in other carbonyls) and can stop; a reader who sees a band at 1100~\wn{} has learned almost nothing and should look elsewhere; a reader who does not see a band at 690~\wn{} can set the monosubstituted ring aside. Two cautions come with this. \LRp{} depends on what else is in the collection, and the two collections show how far it moves: within a factor of two for most rules, but from 43 to 12 for the ester when anhydrides and lactones enter the picture. And \LRn{} depends on which members of a group the collection holds, and travels less well than \LRp{}; the rule-out values should be taken as indicative until they are replicated on further collections.

The study has limits that the reader should weigh. The rules come from one source; other tables give other windows, and the audit says nothing about them, although the method applies to any of them unchanged. Ground truth is the deposited structure matched by our SMARTS definitions, which encode one reading of the source's group names; a different reading of ``saturated'' or ``aromatic'' would move some rules. Band detection is a fixed algorithm with a fixed threshold, and the height measure it uses penalises bands that are strong in area but low in height, the acid O-H envelope above all. Chemotion is an ATR collection with a 93\% aromatic, synthesis-laboratory composition; NIST pools KBr, mull, solution and neat spectra from several decades of instruments, and 43\% of its primary spectra are digitised micrometre scans; we did not stratify by preparation, and mulls in particular add hydrocarbon bands that inflate the methyl and methylene false positives. Each rule was scored alone, without the context a human reader uses. No correction for multiple comparisons was made, and none is claimed. What survives these limits is the pattern: two collections that share none of them agree on it.

Three extensions follow. Other rule sources, and in particular the tables of the classic monographs, can be transcribed and scored with the same scripts against the same collections, which would show how much the sources agree once they are measured rather than compared by eye. Combinations of rules, which is how the fingerprint is actually read, can be scored the same way, with the frozen set as the starting point. And any laboratory with a collection of spectra and structures can compute its own column of the table, which is the only way the composition dependence of the ratios can be mapped. The rule table, the definitions and the scripts are released so that this can be done without re-deriving them.
